\ssr{ВВЕДЕНИЕ}

В последние десятилетия методы компьютерной графики получили широкое распространение в самых разных областях: от индустрии развлечений и компьютерных игр до инженерных расчётов, имитационного моделирования и систем визуализации научных данных. Существенную роль в формировании визуального восприятия сцены играют реалистичные специальные эффекты -- такие, как дым, огонь, туман, взрывы и другие объекты сложной формы. Корректная визуализация подобных явлений позволяет повысить наглядность и выразительность трёхмерных сцен, однако их моделирование связано с рядом вычислительных и алгоритмических сложностей.

Одним из наиболее распространённых подходов к отображению огня и других подобъектов является метод систем частиц. В рамках этого подхода сложный объект представляется большим количеством элементарных частиц, каждая из которых обладает собственной позицией, скоростью, временем жизни и набором визуальных параметров. Совокупное поведение частиц позволяет воспроизвести характерную динамику и структуру пламени. 

В рамках данной курсовой работы рассматривается задача визуализации трёхмерного огня методом систем частиц с учётом параметров внешнего освещения. Объектом исследования в работе является программная система визуализации трёхмерной сцены с анимированными эффектами огня. Предметом исследования являются методы визуализации огня на основе систем частиц и упрощённые модели внешнего освещения, применяемые при рендеринге частиц.

Целью курсовой работы является разработка программного обеспечения для визуализации трёхмерного огня методом систем частиц с учётом внешнего освещения и в исследовании влияния количества частиц и параметров системы на производительность.

Для достижения поставленной цели необходимо выполнить следующие задачи:

\begin{enumerate}
	\item провести сравнение существующих методов визуализации систем частиц и моделей освещения, применяемых для отображения огня и схожих по природе явлений, с точки зрения визуального качества и вычислительной сложности. Выбрать и обосновать алгоритмы визуализации частиц и модель освещения, обеспечивающие требуемый баланс между производительностью и реалистичностью изображения.
	\item разработать структуру программного обеспечения, обеспечивающие представление частиц, систем частиц, камеры, источников света и сцены в целом.
	\item реализовать программный комплекс, позволяющий пользователю создавать и настраивать несколько независимых систем частиц (скорость, размер, цвет, количество частиц), а также управлять параметрами внешнего освещения и виртуальной камеры в реальном времени.
	\item провести исследование зависимости частоты кадров от общего количества частиц в сцене и выбранных алгоритмов визуализации, проанализировать полученные результаты и сформулировать выводы о влиянии выбранных алгоритмов и настроек на производительность.
\end{enumerate}

\clearpage
